\begin{figure}[htbp]
\centering
\begin{tikzpicture}[causal,x=1cm,y=1cm,every node/.style={font=\small,inner sep=2pt},
 neuron/.style={circle,draw=blue!70!black,fill=blue!45,minimum size=2.5mm,inner sep=0}]
\draw[green!55!black,dashed] (0,-2)--(0,1.3)--(1.3,2.3)--(1.3,-1)--cycle;
\draw[red!75] (.7,-1.5)--(.7,1.6);
\draw[purple] (2.6,-1.8)--(2.6,1.1)--(3.7,2.1)--(3.7,-.8)--cycle;
\draw[purple] (4.8,-1.8)--(4.8,1.1)--(5.9,2.1)--(5.9,-.8)--cycle;
\node[draw,anchor=east] (c314input) at (0,.15) {输入};
\draw[->,green!55!black] (c314input.east)--(.7,.15);
\draw[->,magenta] (.7,.15)--(3.1,.15);
\draw[->,red!75] (3.1,.15)--(5.6,.15);
\node at (3.1,2.6) {卷积层};
\node at (5.3,2.6) {卷积层};
\node at (8.2,2.6) {全连接层};
\foreach \i/\yy in {1/1.5,2/.65,3/-.2,4/-1.05}
 \node[neuron] (c314a\i) at (5.6,\yy) {};
\foreach \i/\yy in {1/1.9,2/.45,3/-1}
 \node[neuron] (c314b\i) at (7.7,\yy) {};
\foreach \i/\yy in {1/1.15,2/-.55}
 \node[neuron] (c314c\i) at (9.8,\yy) {};
\foreach \i in {1,2,3,4}\foreach \j in {1,2,3}
 \draw[->,magenta!75] (c314a\i)--(c314b\j);
\foreach \i in {1,2,3}\foreach \j in {1,2}
 \draw[->,magenta!75] (c314b\i)--(c314c\j);
\node[anchor=west] (c314out) at (11.1,.3) {输出预测};
\draw[->,blue!65] (c314c1)--(c314out.west);
\draw[->,blue!65] (c314c2)--(c314out.west);
\end{tikzpicture}
\origfig{14}{卷积神经网络的架构。}
\label{fig:ch03-14}
\end{figure}
